\section*{4. Generosity Indices}
\label{sec:indices}

\subsection*{Index construction methodology}

The Stage 2 pipeline constructs standardized generosity indices across 13 core topics (12 non-salary domains plus wages), following the methodology in \texttt{indices/METHODOLOGY.md}. Each numeric provision is signed (positive indicating higher generosity towards workers), pooled across the entire contract corpus, winsorized to contain extreme leverage points, and standardized into a z-score relative to the empirical sample distribution.

For each topic, two distinct tracks are measured:
\begin{enumerate}
    \item \textbf{Magnitude track}: standardized intensity of provisions where present (e.g., wage levels, vacation hours, overtime surcharge percentages, sick leave wage replacement rates).
    \item \textbf{Coverage track}: the breadth of contractual protections (e.g., the proportion of potential fringe benefits, wellbeing protocols, or leave types provided).
\end{enumerate}
The topic-level combined score represents the unweighted mean of its standardized magnitude and coverage tracks. The overall agreement-level generosity index (\texttt{overall\_z}) is computed as an equal-weighted composite across all eligible topic domains. In addition, an institutional statutory benchmark (``pseudo-CAO'') is continuously tracked, scoring mandatory legal statutory baselines under Dutch labor law (e.g., WML, statutory holiday entitlements, WAZO parental leave, and statutory transition compensation).

\subsection*{Overall generosity trend and statutory baseline}

Figure~\ref{fig:indices_overall_trend} plots the evolution of the overall \emph{combined} generosity score (\texttt{overall\_z}, magnitude and coverage tracks together) across the monthly agreement panel from \IndPanelStartMonth\ to \IndPanelEndMonth\ (\IndPanelTotalMonths\ months). ``Mature'' months are defined as those with at least \IndMatureMinCaos\ active CAOs, which excludes only the earliest, thinly-covered years of the corpus.

\begin{figure}[H]
    \centering
    \includegraphics[width=\textwidth]{figures/indices/indices_overall_z_trend.png}
    \caption{Overall CAO generosity trend and statutory benchmark over time (monthly panel).}
    \label{fig:indices_overall_trend}
\end{figure}

The mean overall generosity score tracks steadily above the statutory floor throughout the sample period. While the statutory benchmark (red dashed line) shifts upwards with statutory reforms (such as minimum wage indexation, statutory vacation leave expansions, and the introduction of paid parental leave schemes), CAO provisions consistently maintain a positive generosity margin over the legal minimum. The 10th--90th percentile band highlights sustained heterogeneity across agreements throughout the sample period.

\subsection*{Domain generosity distributions}

Figure~\ref{fig:indices_topic_distributions} displays boxplots for the cross-sectional distribution of standardized generosity scores across individual domains. Table~\ref{tab:indices_topline} provides corresponding summary statistics.

\begin{figure}[H]
    \centering
    \includegraphics[width=\textwidth]{figures/indices/indices_topic_distributions.png}
    \caption{Distribution of domain generosity scores across collective agreements.}
    \label{fig:indices_topic_distributions}
\end{figure}

\begin{table}[H]
    \centering
    \caption{Topline summary statistics for overall and domain generosity indices}
    \label{tab:indices_topline}
    \begin{tabular}{lrrrrr}
    \hline\hline
    Domain / Index & Mean & Std & P10 & Median & P90 \\
    \hline
    Overall Generosity (Combined) & -0.01 & 0.38 & -0.46 & -0.02 & 0.45 \\
    Leave Provisions & -0.00 & 0.66 & -0.91 & -0.06 & 0.98 \\
    Sickness \& Care & -0.00 & 0.57 & -0.74 & 0.06 & 0.76 \\
    Pensions & -0.04 & 0.77 & -0.80 & 0.01 & 0.91 \\
    Overtime & 0.00 & 0.56 & -1.22 & 0.17 & 0.41 \\
    Bonuses & -0.08 & 0.78 & -0.94 & -0.21 & 1.17 \\
    Fringe Benefits & 0.07 & 0.80 & -0.92 & -0.07 & 1.14 \\
    Training & 0.02 & 0.69 & -0.61 & 0.25 & 0.66 \\
    Contract Types & 0.00 & 0.58 & -0.69 & 0.01 & 0.78 \\
    Termination & -0.01 & 0.54 & -0.88 & 0.08 & 0.70 \\
    Home Office & 0.01 & 0.96 & -0.47 & -0.39 & 1.64 \\
    \hline\hline
\end{tabular}

\end{table}

Variation is especially pronounced in \textbf{Home Office} (standard deviation \IndSDHomeoffice) and \textbf{Fringe Benefits} (\IndSDFringe), where agreements diverge substantially between comprehensive packages and minimal statutory compliance. By contrast, \textbf{Termination} (\IndSDTermination) and \textbf{Overtime} (\IndSDOvertime) exhibit tighter clustering due to strict national legal framing under the Civil Code (BW) and Working Hours Act (ATW).

\subsection*{The wage ladder vs. statutory minimum wage}

Figure~\ref{fig:indices_wage_ladder} documents the evolution of the contractual wage ladder relative to the statutory minimum wage (WML). For each CAO and calendar year, the 10th percentile (\texttt{mw\_low}), median (\texttt{mw\_median}), and 90th percentile (\texttt{mw\_high}) of the full-time wage scale are evaluated against the statutory monthly minimum wage. Years with fewer than \IndWageYearMinRows\ wage-scale observations are excluded to avoid noisy medians from a thin tail.

\begin{figure}[H]
    \centering
    \includegraphics[width=\textwidth]{figures/indices/indices_wage_ladder_vs_wml.png}
    \caption{CAO wage scale ladder vs.\ statutory minimum wage (WML), \IndWageYearMin--\IndWageYearMax.}
    \label{fig:indices_wage_ladder}
\end{figure}

In recent cohorts, the median CAO scale wage reached \IndWageMedianLatest, corresponding to approximately \IndWageToWmlRatio\ times the statutory floor (\IndWmlMonthLatest). The scale floor (\texttt{mw\_low}) tracks closely above the WML line, demonstrating immediate pass-through of statutory minimum wage updates into lower-tier collective bargaining scales.

\subsection*{Top and bottom agreements}

Table~\ref{tab:indices_top_bottom} ranks the highest and lowest collective agreements by composite generosity in their latest observable, quality-gated contracts. Because a document's \texttt{overall\_z} is an available-case mean, one scored on very few topics can register a noisy extreme value that looks like genuine (un)generosity rather than a measurement artifact; the ranking therefore excludes carried/thin documents and any document scoring on fewer than \IndMinTopicsScored\ of the available topics (\IndRankingExcludedDocs\ of \IndCompositeTotalDocs\ documents excluded on this basis), leaving \IndRankingRankableCaos\ of \IndRankingTotalCaos\ CAOs rankable.

\begin{table}[H]
    \centering
    \caption{Top and bottom collective agreements by composite generosity (\texttt{overall\_z})}
    \label{tab:indices_top_bottom}
    \begin{tabular}{llp{6.5cm}rr}
    \hline\hline
    Rank & CAO & Document Name & Overall z & Percentile \\
    \hline
    Top 1 & 1646 & CAO Rijk 2024-2025 per 1 juli 2025 - schoon & 1.34 & 89.3\% \\
    Top 2 & 1022 & CAO MBO 2024 & 1.29 & 89.6\% \\
    Top 3 & 156 & Cao Ziekenhuizen 2023-2025 & 1.15 & 86.9\% \\
    Top 4 & 625 & CAO 2024 2025def & 1.07 & 88.9\% \\
    Top 5 & 317 & Tekst CAO Gehandicaptenzorg 2025-2026 (ID 161822) & 1.04 & 84.4\% \\
    Bottom 5 & 1429 & RTO cao 2025tekenversiewordWAV & -1.28 & 37.1\% \\
    Bottom 4 & 4000 & A.08032024 cao PAWW Bouw 3no 02A & -1.11 & 41.7\% \\
    Bottom 3 & 4171 & CAO onwerkbaar weer BIKUDAK 2023 - 2024 ttw dec… & -1.05 & 43.2\% \\
    Bottom 2 & 558 & 20130222 CAO Papiergroothandel 2010 2013 & -1.01 & 41.3\% \\
    Bottom 1 & 4062 & 03 - Cao DISS 2024 tm 2028 ongeren na 3e TTW & -0.96 & 45.3\% \\
    \hline\hline
\end{tabular}

\end{table}

The top-ranked agreements (e.g., CAO Rijk, CAO MBO, and CAO Ziekenhuizen) feature extensive leave entitlements, high employer pension contributions, robust training funding, and structured wage ladders. The bottom-ranked documents remain primarily single-purpose agreements, bad-weather compensation schemes (onwerkbaar weer), or transitional supplementary protocols; their domain coverage clears the gate above but is still narrow, consistent with genuinely limited bargaining scope rather than a data-quality artifact.

\subsection*{Statutory-response elasticities}

Table~\ref{tab:indices_statutory_response} reports $\beta$ on contemporaneous statutory change (\texttt{d\_stat\_t}) from the statutory-response regressions in \texttt{indices/ADVANCED\_ANALYSIS.md} (year-clustered SEs). A $\beta$ near 1.0 indicates CAO provisions move roughly one-for-one with statutory reforms; leave provisions do so closely, while the aggregate overall series responds more strongly still, reflecting reforms that shift several domains at once in the same statutory event.

\begin{table}[H]
    \centering
    \caption{Statutory-response elasticities by domain}
    \label{tab:indices_statutory_response}
    \begin{tabular}{lrrrr}
    \hline\hline
    Domain & $\beta$ (d\_stat\_t) & Clustered SE & $t$ & N \\
    \hline
    Leave Provisions & 1.08 & 0.31 & 3.49 & 3,218 \\
    Contract Types & 0.28 & 0.03 & 8.67 & 3,218 \\
    Overall & 1.14 & 0.54 & 2.10 & 3,218 \\
    \hline\hline
\end{tabular}

\end{table}

\subsection*{Generosity vs.\ pay level}

Table~\ref{tab:indices_pay_correlations} reports, per domain, the correlation between a CAO's combined generosity score and its wage level. Correlations are positive but modest for most domains, and slightly negative for a few (e.g.\ overtime, contract type) -- evidence that generosity and pay are largely independent margins rather than a single underlying quality hierarchy.

\begin{table}[H]
    \centering
    \caption{Correlation between domain generosity and wage level}
    \label{tab:indices_pay_correlations}
    \begin{tabular}{lrrr}
    \hline\hline
    Domain & N & Pearson $r$ & Spearman $\rho$ \\
    \hline
    Leave Provisions & 2,444 & 0.33 & 0.39 \\
    Sickness \& Care & 2,444 & 0.16 & 0.25 \\
    Pensions & 2,444 & 0.12 & 0.17 \\
    Overtime & 2,444 & -0.07 & 0.02 \\
    Bonuses & 2,444 & 0.36 & 0.34 \\
    Fringe Benefits & 2,444 & 0.34 & 0.40 \\
    Training & 2,444 & -0.01 & 0.06 \\
    Contract Types & 2,444 & 0.03 & 0.05 \\
    Termination & 2,444 & 0.19 & 0.20 \\
    Home Office & 2,444 & 0.32 & 0.28 \\
    \hline\hline
\end{tabular}

\end{table}

\subsection*{Multivariate structure and methodological caveats}

Table~\ref{tab:indices_factor_loadings} reports the 3-factor loadings on the combined (magnitude+coverage) domain scores from the factor analysis in \texttt{indices/ADVANCED\_ANALYSIS.md}. No single factor loads highly across all domains, confirming that collective bargaining generosity does not collapse into one unidimensional axis. The overall Kaiser-Meyer-Olkin (KMO) measure of sampling adequacy is \IndKmoStat, indicating moderate common variance across topics.

\begin{table}[H]
    \centering
    \caption{Three-factor loadings on combined domain generosity scores}
    \label{tab:indices_factor_loadings}
    \begin{tabular}{lrrrr}
    \hline\hline
    Domain & F1 & F2 & F3 & Communality \\
    \hline
    Wage Scale & 0.00 & 0.70 & 0.12 & 0.51 \\
    Leave Provisions & 0.16 & 0.28 & 0.77 & 0.69 \\
    Sickness \& Care & 0.61 & 0.13 & 0.14 & 0.41 \\
    Pensions & 0.44 & 0.14 & 0.06 & 0.21 \\
    Overtime & 0.35 & -0.11 & 0.34 & 0.25 \\
    Bonuses & 0.38 & 0.46 & 0.14 & 0.37 \\
    Fringe Benefits & 0.40 & 0.46 & 0.04 & 0.38 \\
    Training & 0.48 & 0.00 & 0.07 & 0.24 \\
    Contract Types & 0.47 & 0.04 & 0.15 & 0.24 \\
    Termination & 0.52 & 0.26 & -0.03 & 0.35 \\
    Home Office & 0.02 & 0.41 & 0.40 & 0.33 \\
    \hline\hline
\end{tabular}

\end{table}

Several methodological caveats apply to the indices:
\begin{itemize}
    \item \textbf{Pension treatment}: pension generosity is scored on an \emph{available-case only} basis. Blank entries reflect fund-level administration and sectoral deferrals rather than contractual absence, and are never zero-filled.
    \item \textbf{Open judgment items}: \IndOpenJudgmentCells\ edge-case cells remain designated as open judgment items pending expert adjudication (\texttt{indices/review/all\_open\_cant\_tells.csv}).
    \item \textbf{Statutory fingerprinting}: \IndFingerprintSuspects\ provisions identified by automated screening as identical to statutory defaults have not yet been individually adjudicated (\texttt{indices/out/statutory\_fingerprint\_suspects.csv}).
    \item \textbf{AI exclusion}: Because explicit artificial intelligence clauses are present in only \NonSalAiShare\ of CAOs, AI provisions are monitored descriptively but excluded from the headline composite generosity metric (\texttt{overall\_z}).
\end{itemize}

